%! AP Statistics — Unit 2, Lesson 2.2 — Group Activity
\documentclass[10pt]{article}
\usepackage{apstats-article}
\usepackage{apstats-boxes}

%=============================================================================
\begin{document}

\pageheader{Unit 2, Lesson 2.2}{Group Activity}
\namepartnerperiod

\vspace{0.1in}

\begin{tcolorbox}[colback=sky, colframe=navy, boxrule=0.8pt, arc=2mm,
    left=3mm, right=3mm, top=2mm, bottom=2mm]
\small
Work with your group on the tier assigned by your teacher.
For each table: (1) compute conditional relative frequencies, (2) decide whether
the two variables are \textbf{associated}, and (3) justify using specific values.
\end{tcolorbox}

\vspace{0.12in}

% ════════════════════════════════════════════════════════════════════════════
% TIER R
% ════════════════════════════════════════════════════════════════════════════
\begin{tcolorbox}[colback=navylight!30, colframe=navy, boxrule=0.9pt, arc=2mm,
    left=3mm, right=3mm, top=2mm, bottom=2mm,
    title={\bfseries Tier R --- Remediate}, fonttitle=\bfseries\small, halign title=center]
\small

\textbf{Context:} 120 students were surveyed on two variables: whether they have an
after-school job (Job: Yes/No) and whether they play a school sport (Sport: Yes/No).

\vspace{8pt}
\renewcommand{\arraystretch}{1.8}
\begin{tabularx}{\linewidth}{@{} l C{2.4cm} C{2.4cm} C{2.2cm} @{}}
\rowcolor{sky}
             & \textbf{Sport: Yes} & \textbf{Sport: No} & \textbf{Row Total} \\
\hline
\textbf{Job: Yes} & 12 & 28 & 40 \\
\textbf{Job: No}  & 48 & 32 & 80 \\
\hline
\textbf{Column Total} & 60 & 60 & 120 \\
\end{tabularx}

\vspace{10pt}
\textbf{1.} Compute the conditional relative frequencies for each \emph{job group}
(divide each row by its row total):

\vspace{6pt}
\renewcommand{\arraystretch}{1.8}
\begin{tabularx}{\linewidth}{@{} l C{2.6cm} C{2.6cm} C{1.8cm} @{}}
\rowcolor{sky}
             & \textbf{Sport: Yes} & \textbf{Sport: No} & \textbf{Total} \\
\hline
Job: Yes & \blank{2.2cm} & \blank{2.2cm} & 100\% \\
Job: No  & \blank{2.2cm} & \blank{2.2cm} & 100\% \\
\end{tabularx}

\vspace{10pt}
\textbf{2.} Is there evidence of an association between having a job and playing a sport?
Circle: \quad \textbf{Yes} \quad / \quad \textbf{No}

\vspace{6pt}
Justify using two specific conditional percentages:\\[6pt]
\writeline\\[1pt]\writeline

\end{tcolorbox}

\vspace{0.14in}

% ════════════════════════════════════════════════════════════════════════════
% TIER A
% ════════════════════════════════════════════════════════════════════════════
\begin{tcolorbox}[colback=greenbg, colframe=greenacc, boxrule=0.9pt, arc=2mm,
    left=3mm, right=3mm, top=2mm, bottom=2mm,
    title={\bfseries Tier A --- Approaching Proficiency}, fonttitle=\bfseries\small, halign title=center]
\small

\textbf{Context:} 180 students were asked their class year (Freshman / Sophomore) and their
preferred movie genre (Action / Comedy / Drama). Here are the counts:

\vspace{6pt}
\renewcommand{\arraystretch}{1.8}
\begin{tabularx}{\linewidth}{@{} l C{2.2cm} C{2.2cm} C{2.2cm} C{2.2cm} @{}}
\rowcolor{sky}
                    & \textbf{Action} & \textbf{Comedy} & \textbf{Drama} & \textbf{Row Total} \\
\hline
\textbf{Freshman}   & 36 & 54 & 30 & \blank{1.8cm} \\
\textbf{Sophomore}  & 30 & 18 & 12 & \blank{1.8cm} \\
\hline
\textbf{Col.\ Total}& \blank{1.6cm} & \blank{1.6cm} & \blank{1.6cm} & \blank{1.6cm} \\
\end{tabularx}

\vspace{10pt}
\textbf{1.} Fill in all marginal totals. Then compute the conditional relative
frequencies for each class year:

\vspace{6pt}
\renewcommand{\arraystretch}{1.8}
\begin{tabularx}{\linewidth}{@{} l C{2.4cm} C{2.4cm} C{2.4cm} C{1.6cm} @{}}
\rowcolor{sky}
                    & \textbf{Action} & \textbf{Comedy} & \textbf{Drama} & \textbf{Total} \\
\hline
Freshman   & \blank{2cm} & \blank{2cm} & \blank{2cm} & 100\% \\
Sophomore  & \blank{2cm} & \blank{2cm} & \blank{2cm} & 100\% \\
\end{tabularx}

\vspace{10pt}
\textbf{2.} In the space below, sketch a \textbf{segmented bar chart} for these
conditional distributions. Label the axes and use a legend.

\vspace{0.7in}

\textbf{3.} Is there evidence of an association between class year and genre preference?
State your conclusion and justify with specific values.\\[6pt]
\writeline\\[1pt]\writeline\\[1pt]\writeline

\end{tcolorbox}

\vspace{0.14in}

% ════════════════════════════════════════════════════════════════════════════
% TIER E
% ════════════════════════════════════════════════════════════════════════════
\begin{tcolorbox}[colback=redbg, colframe=redacc, boxrule=0.9pt, arc=2mm,
    left=3mm, right=3mm, top=2mm, bottom=2mm,
    title={\bfseries Tier E --- Extension}, fonttitle=\bfseries\small, halign title=center]
\small

\textbf{Context:} 300 adults were asked about their commute method (Car / Transit / Walk) and
their reported stress level (Low / High). Here are the \emph{joint relative frequencies}:

\vspace{6pt}
\renewcommand{\arraystretch}{1.8}
\begin{tabularx}{\linewidth}{@{} l C{2.4cm} C{2.4cm} C{2.4cm} @{}}
\rowcolor{sky}
                  & \textbf{Stress: Low} & \textbf{Stress: High} & \textbf{Row Total} \\
\hline
\textbf{Car}     & 0.20 & 0.30 & \blank{1.8cm} \\
\textbf{Transit} & 0.15 & 0.15 & \blank{1.8cm} \\
\textbf{Walk}    & 0.14 & 0.06 & \blank{1.8cm} \\
\hline
\textbf{Col.\ Total} & \blank{1.6cm} & \blank{1.6cm} & \blank{1.6cm} \\
\end{tabularx}

\vspace{10pt}
\textbf{1.} Verify: do the joint relative frequencies in the table sum to 1.00?
Show your check.\\[6pt]
\writeline

\vspace{8pt}
\textbf{2.} Compute the \textbf{conditional} relative frequency of high stress
\emph{given each commute method}. (Hint: divide each row's ``High'' value by its row total.)

\vspace{6pt}
\renewcommand{\arraystretch}{1.8}
\begin{tabularx}{\linewidth}{@{} l C{3.4cm} @{}}
\rowcolor{sky} \textbf{Commute} & \textbf{P(High stress $|$ commute)} \\
\hline
Car     & \blank{2.8cm} \\
Transit & \blank{2.8cm} \\
Walk    & \blank{2.8cm} \\
\end{tabularx}

\vspace{8pt}
\textbf{3.} Write a two-sentence conclusion: Is there evidence of an association between
commute method and stress level? Which group has the highest conditional stress rate?\\[6pt]
\writeline\\[1pt]\writeline\\[1pt]\writeline

\vspace{6pt}
\textbf{4. Stretch:} If you were to condition on stress level instead (compute the
conditional distribution of commute method \emph{given} stress level), would the
association conclusion change? Try it and explain.\\[6pt]
\writeline\\[1pt]\writeline

\end{tcolorbox}

\end{document}
